\scanpage{59}
Si $M$ est un motif sur $K$, $1/2$ simple pour simplifier, alors on a
\[
\unsure{\mathfrak g|g}_\ell(\overset{\unsure{\ell|s}}{\bar K}/K)\longrightarrow\operatorname{End}(M_l)\simeq\check M_l\otimes M_l
\]
On pense que l'image provient d'un motif $\mathfrak g_{\deleted{(M)}M}\inserted{{}\subset\check M\otimes M}$ indépendant de $l$, [qui est un motif en \unsure{Algèbre|algèbre|algèbres} de Lie].

Pour $M$ de plus en plus grand, on trouve un \uline{pro-motif} \inserted{(en alg de Lie)} canonique, appelé \uline{promotif fondamental}, qui peut aussi se construire de façon canonique en termes de la catégorie des motifs, munie de ses structures tensorielles \ldots

Quelle est la partie ``algébrique'' de ce promotif ? [Sur un corps fini, c'est tout -- le promotif n'étant lui-$=$ autre que $\mathbb Z(0)$ ! ; \deleted{dans}\inserted{en} tout \deleted{\unsure{autre}}\inserted{autre} cas, \deleted{\unsure{on}} \inserted{je pense qu'il est \uline{nul} \ldots.} \deleted{\unsure{il s'adresse|il s'intéresse} contenu dans le \uline{centre} du dit promotif fondamental, et je pense que c'est réduit à $\mathbb Z(0)$ [\unsure{n-} ``scalaires'']}
% archive sheet number: 58

